% GENERATED FILE. Edit canonical lessons, then run npm run generate:books.
% canonical-source-hash: fnv1a64:68588c552772badb
% canonical-lessons: ML-C422-niyamasabha, ML-C422-amgam, ML-C422-pratinidhi, ML-C422-pratinidhikarikkuka, ML-C422-pratipaksam

\chapter{The Assembly}
\label{ch:ml-the-assembly}
\hlchaptermodality{422}

\begin{chapteropening}
\textbf{By the end of this chapter:} \emph{I can say an assembly, a member, a representative, to represent and the opposition.}

Complete the last lesson of chapter 422: I can say an assembly, a member, a representative, to represent and the opposition.
\end{chapteropening}

\section[\texorpdfstring{\ml{നിയമസഭ} (niyamasabha)}{നിയമസഭ (niyamasabha)}]{\ml{നിയമസഭ} (niyamasabha) — a legislative assembly}

\label{lesson:ML-C422-niyamasabha}



\begin{quote}
Before the new one: say the Malayalam for a chief minister, then the Malayalam for a governor.
\end{quote}

\subsection*{You'll want to know: \ml{നിയമസഭ}}
\textbf{\ml{നിയമസഭ}} — \emph{niyamasabha} — \textquotedblleft{}a legislative assembly\textquotedblright{}.

\begin{quote}
\ml{നിയമസഭയിൽ} \ml{ഇന്ന്} \ml{വലിയ} \ml{ചർച്ച} \ml{നടന്നു}.
\end{quote}

\begin{quote}
\emph{niyamasabhayil innŭ valiya carcca naṭannu.}
\end{quote}

\begin{quote}
There was a big discussion in the assembly today.
\end{quote}

A member of it is \ml{നിയമസഭാംഗം}.

\begin{cousinweb}
\textbf{\ml{നിയമസഭ}} (\emph{niyamasabha}) is a Malayalam compound built from Sanskrit \emph{niyama}, \textquotedblleft{}a fixed law\textquotedblright{} (CDIAL 7250), and \emph{sabhā}, \textquotedblleft{}an assembly, a hall\textquotedblright{} (CDIAL 13167): the assembly that makes laws.
\end{cousinweb}

\begin{grammarlens}[title={What you've built}]
One more word for giving reasons and opinions.
\end{grammarlens}

\subsection*{Your turn}
\begin{itemize}
\raggedright
  \item \emph{Say it:} \emph{niyamasabha}
  \item \emph{Say it:} \emph{niyamasabha}, once more
  \item \emph{Say it:} \emph{niyamasabhayil innŭ valiya carcca naṭannu.} — There was a big discussion in the assembly today.
  \item \emph{Recall:} say the Malayalam for a chief minister, then the Malayalam for a governor, then say \emph{niyamasabha} again
\end{itemize}

\begin{tcolorbox}[breakable,colback=teal!4,colframe=teal!35!black,title={Before you move on}]
What does \ml{നിയമസഭ} mean? (\textquotedblleft{}A legislative assembly\textquotedblright{}.) Say it once more.
\end{tcolorbox}
\section[\texorpdfstring{\ml{അംഗം} (aṁgaṁ)}{അംഗം (aṁgaṁ)}]{\ml{അംഗം} (aṁgaṁ) — a member}

\label{lesson:ML-C422-amgam}



\begin{quote}
Before the new one: say the Malayalam for a governor, then the Malayalam for a legislative assembly.
\end{quote}

\subsection*{You'll want to know: \ml{അംഗം}}
\textbf{\ml{അംഗം}} — \emph{aṁgaṁ} — \textquotedblleft{}a member\textquotedblright{}.

\begin{quote}
\ml{ഞാൻ} \ml{ആ} \ml{ക്ലബ്ബിലെ} \ml{അംഗമാണ്}.
\end{quote}

\begin{quote}
\emph{ñān ā klabbile aṁgamāṇŭ.}
\end{quote}

\begin{quote}
I am a member of that club.
\end{quote}

Plural: \ml{അംഗങ്ങൾ}, \emph{aṁgaṅṅaḷ}.

Plural: \ml{അംഗങ്ങൾ}.

\begin{cousinweb}
\textbf{\ml{അംഗം}} (\emph{aṁgaṁ}) is borrowed from Sanskrit \emph{aṅga}, \textquotedblleft{}a limb; the body\textquotedblright{} (CDIAL 114). A member is a limb of a body of people, as in English.
\end{cousinweb}

\begin{grammarlens}[title={What you've built}]
One more word for giving reasons and opinions.
\end{grammarlens}

\subsection*{Your turn}
\begin{itemize}
\raggedright
  \item \emph{Say it:} \emph{aṁgaṁ}
  \item \emph{Say it:} \emph{aṁgaṁ}, once more
  \item \emph{Say it:} \emph{ñān ā klabbile aṁgamāṇŭ.} — I am a member of that club.
  \item \emph{Recall:} say the Malayalam for a governor, then the Malayalam for a legislative assembly, then say \emph{aṁgaṁ} again
\end{itemize}

\begin{tcolorbox}[breakable,colback=teal!4,colframe=teal!35!black,title={Before you move on}]
What does \ml{അംഗം} mean? (\textquotedblleft{}A member\textquotedblright{}.) Say it once more.
\end{tcolorbox}
\section[\texorpdfstring{\ml{പ്രതിനിധി} (pratinidhi)}{പ്രതിനിധി (pratinidhi)}]{\ml{പ്രതിനിധി} (pratinidhi) — a representative}

\label{lesson:ML-C422-pratinidhi}



\begin{quote}
Before the new one: say the Malayalam for a legislative assembly, then the Malayalam for a member.
\end{quote}

\subsection*{You'll want to know: \ml{പ്രതിനിധി}}
\textbf{\ml{പ്രതിനിധി}} — \emph{pratinidhi} — \textquotedblleft{}a representative\textquotedblright{}.

\begin{quote}
\ml{ഓരോ} \ml{ക്ലാസിൽ} \ml{നിന്നും} \ml{ഒരു} \ml{പ്രതിനിധി} \ml{യോഗത്തിന്} \ml{വന്നു}.
\end{quote}

\begin{quote}
\emph{ōrō klāsil ninnuṁ oru pratinidhi yōgattinŭ vannu.}
\end{quote}

\begin{quote}
A representative came to the meeting from each class.
\end{quote}

Plural: \ml{പ്രതിനിധികൾ}.

\begin{cousinweb}
\textbf{\ml{പ്രതിനിധി}} (\emph{pratinidhi}) is borrowed from Sanskrit \emph{pratinidhi}, \textquotedblleft{}a substitute, a representative\textquotedblright{} (Monier-Williams): one put in the place of another.
\end{cousinweb}

\begin{grammarlens}[title={What you've built}]
One more word for giving reasons and opinions.
\end{grammarlens}

\subsection*{Your turn}
\begin{itemize}
\raggedright
  \item \emph{Say it:} \emph{pratinidhi}
  \item \emph{Say it:} \emph{pratinidhi}, once more
  \item \emph{Say it:} \emph{ōrō klāsil ninnuṁ oru pratinidhi yōgattinŭ vannu.} — A representative came to the meeting from each class.
  \item \emph{Recall:} say the Malayalam for a legislative assembly, then the Malayalam for a member, then say \emph{pratinidhi} again
\end{itemize}

\begin{tcolorbox}[breakable,colback=teal!4,colframe=teal!35!black,title={Before you move on}]
What does \ml{പ്രതിനിധി} mean? (\textquotedblleft{}A representative\textquotedblright{}.) Say it once more.
\end{tcolorbox}
\section[\texorpdfstring{\ml{പ്രതിനിധീകരിക്കുക} (pratinidhīkarikkuka)}{പ്രതിനിധീകരിക്കുക (pratinidhīkarikkuka)}]{\ml{പ്രതിനിധീകരിക്കുക} (pratinidhīkarikkuka) — to represent}

\label{lesson:ML-C422-pratinidhikarikkuka}



\begin{quote}
Before the new one: say the Malayalam for a member, then the Malayalam for a representative.
\end{quote}

\subsection*{You'll want to know: \ml{പ്രതിനിധീകരിക്കുക}}
\textbf{\ml{പ്രതിനിധീകരിക്കുക}} — \emph{pratinidhīkarikkuka} — \textquotedblleft{}to represent\textquotedblright{}.

\begin{quote}
\ml{അവൾ} \ml{ഞങ്ങളുടെ} \ml{കോളേജിനെ} \ml{പ്രതിനിധീകരിച്ചു}.
\end{quote}

\begin{quote}
\emph{avaḷ ñaṅṅaḷuṭe kōḷējine pratinidhīkariccu.}
\end{quote}

\begin{quote}
She represented our college.
\end{quote}

Past: \ml{പ്രതിനിധീകരിച്ചു}, \emph{pratinidhīkariccu}.

What is represented takes -\ml{നെ}.

\begin{cousinweb}
\textbf{\ml{പ്രതിനിധീകരിക്കുക}} (\emph{pratinidhīkarikkuka}) is built from Sanskrit \emph{pratinidhi}, \textquotedblleft{}a representative\textquotedblright{} (Monier-Williams), with -\ml{ീകരിക്കുക}, \textquotedblleft{}to make into\textquotedblright{}: to be made the representative.
\end{cousinweb}

\begin{grammarlens}[title={What you've built}]
One more word for giving reasons and opinions.
\end{grammarlens}

\subsection*{Your turn}
\begin{itemize}
\raggedright
  \item \emph{Say it:} \emph{pratinidhīkarikkuka}
  \item \emph{Say it:} \emph{pratinidhīkarikkuka}, once more
  \item \emph{Say it:} \emph{avaḷ ñaṅṅaḷuṭe kōḷējine pratinidhīkariccu.} — She represented our college.
  \item \emph{Recall:} say the Malayalam for a member, then the Malayalam for a representative, then say \emph{pratinidhīkarikkuka} again
\end{itemize}

\begin{tcolorbox}[breakable,colback=teal!4,colframe=teal!35!black,title={Before you move on}]
What does \ml{പ്രതിനിധീകരിക്കുക} mean? (\textquotedblleft{}To represent\textquotedblright{}.) Say it once more.
\end{tcolorbox}
\section[\texorpdfstring{\ml{പ്രതിപക്ഷം} (pratipakṣaṁ)}{പ്രതിപക്ഷം (pratipakṣaṁ)}]{\ml{പ്രതിപക്ഷം} (pratipakṣaṁ) — the opposition}

\label{lesson:ML-C422-pratipaksam}



\begin{quote}
Before the new one: say the Malayalam for a representative, then the Malayalam for to represent.
\end{quote}

\subsection*{You'll want to know: \ml{പ്രതിപക്ഷം}}
\textbf{\ml{പ്രതിപക്ഷം}} — \emph{pratipakṣaṁ} — \textquotedblleft{}the opposition\textquotedblright{}.

\begin{quote}
\ml{പ്രതിപക്ഷം} \ml{പുതിയ} \ml{നികുതിയെ} \ml{എതിർത്തു}.
\end{quote}

\begin{quote}
\emph{pratipakṣaṁ putiya nikutiye etirttu.}
\end{quote}

\begin{quote}
The opposition opposed the new tax.
\end{quote}

Its members sit across from the government in the assembly.

\begin{cousinweb}
\textbf{\ml{പ്രതിപക്ഷം}} (\emph{pratipakṣaṁ}) is built from Sanskrit \emph{prati-}, \textquotedblleft{}against, opposite\textquotedblright{}, and \emph{pakṣa}, \textquotedblleft{}a side\textquotedblright{} (CDIAL 7627): the side opposite.
\end{cousinweb}

\begin{grammarlens}[title={What you've built}]
One more word for giving reasons and opinions.
\end{grammarlens}

\subsection*{Your turn}
\begin{itemize}
\raggedright
  \item \emph{Say it:} \emph{pratipakṣaṁ}
  \item \emph{Say it:} \emph{pratipakṣaṁ}, once more
  \item \emph{Say it:} \emph{pratipakṣaṁ putiya nikutiye etirttu.} — The opposition opposed the new tax.
  \item \emph{Recall:} say the Malayalam for a representative, then the Malayalam for to represent, then say \emph{pratipakṣaṁ} again
\end{itemize}

\begin{tcolorbox}[breakable,colback=teal!4,colframe=teal!35!black,title={Before you move on}]
What does \ml{പ്രതിപക്ഷം} mean? (\textquotedblleft{}The opposition\textquotedblright{}.) Say it once more.
\end{tcolorbox}
